\chapter*{CONCLUSION}
\addcontentsline{toc}{chapter}{CONCLUSION}
\label{ch:conclusion}

This thesis benchmarked six ICS anomaly detectors --- Isolation Forest, One-Class SVM, Dense Autoencoder, LSTM Autoencoder, USAD, and TranAD --- on HAI~21.03~\cite{shin2020hai} and the Morris gas-pipeline corpus~\cite{morris2014industrial} under three time-aware metrics, evaluated cross-testbed transfer under two threshold-calibration regimes with a within-HAI leave-one-process-out control, and evaluated per-sensor attribution at the process level using a multi-seed sensitivity check.  Every figure and every numerical result is reproducible from a committed YAML configuration via the \file{experiments/} command-line interface.

Three findings define the work.  First, source-calibrated cross-testbed transfer collapses every detector to the predict-all trivial baseline of the target's class prior --- F\textsubscript{1}~$\approx 0.95$ on Morris (with a windowed attack rate of $91\%$) and F\textsubscript{1}~$\approx 0.05$ on HAI ($2.4\%$); only target-calibrated transfer reveals a meaningful classical-versus-SOTA separation, with the windowed SOTA models reaching $0.38$--$0.44$ eTaPR on HAI~$\rightarrow$~Morris while the classical baselines do not exceed $0.13$.  Second, within-HAI LOPO shows that feature-distribution shift produces \emph{targeted} detection blindness rather than global degradation, with a within-testbed ceiling of approximately $0.57$ pointwise F\textsubscript{1} that bounds what any schema-level alignment can plausibly achieve.  Third, the weakest detector by F\textsubscript{1} (Dense Autoencoder) is the only AE-family or windowed model whose per-feature attribution beats a random baseline on all three attacked HAI processes, and SHAP-based attribution on the two classical baselines (Isolation Forest, One-Class SVM) places them above-random as well, exhibiting the same detection-versus-localization inversion: the most actionable per-feature attributions come from models near the bottom of the detection table.

The findings translate into three operational recommendations.  Calibrate the decision threshold on target-domain normal data, not on source-domain validation data, before activating a transferred detector.  Report multiple time-aware metrics --- point-wise, point-adjust (with explicit acknowledgement of its inflation~\cite{kim2022towards}), and eTaPR~\cite{hwang2022etapr} --- and treat single-metric rankings with suspicion.  Audit attribution alongside detection where sensor-level localization matters operationally.  A hybrid pipeline that uses a SOTA model for flagging and a Dense Autoencoder or SHAP-explained classical model for attribution may outperform either alone on operator-facing triage.

Beyond the empirical findings, the thesis contributes a reproducible benchmarking protocol and an open-source implementation that future work can extend without re-deriving the apparatus from scratch.  Adding a seventh detector amounts to writing one Python class implementing the abstract \code{AnomalyDetector} base; adding a third dataset amounts to writing one loader function returning the unified schema.  The single methodological recommendation this thesis makes to the ICS anomaly-detection community is to report the full grid --- multiple metrics, multiple calibration regimes, attribution alongside detection --- rather than the single headline F\textsubscript{1} that currently dominates publications.
